\subsection{男性的性反应}

按照玛斯特斯和约翰逊夫妇的理论，男性与女性一样，要经历兴奋期、平台期、高潮期、消退期这四个阶段。所不同的是，男性的变化以阴茎、阴囊和内生殖器官为中心，并波及全身的肌肉与血液循环。在性欲有了最初的萌发后，男性的身体会发生如下变化。

\subsubsection{兴奋期}

男性进入性兴奋期最直接的标志，就是阴茎勃起——阴茎内的海绵体在动脉血流大量涌入时充血胀大，回流的道路被压窄，于是阴茎从松弛变得坚挺。勃起可以在性意念、对生殖器的直接触摸，或是其他感觉刺激下发生，整个过程通常只需要几秒钟。

进入兴奋期后，肾上腺释放肾上腺素，使得心跳加快、血压升高、呼吸急促。阴囊的皮肤收紧、变厚，提睾肌收缩，把睾丸向上提起、贴近身体。有些男性的皮肤上同样会出现红晕。

\subsubsection{平台期}

随着性刺激的继续，阴茎进一步充血，达到完全勃起，龟头的颜色加深，变成深红色或紫色。尿道口会有少量清亮的分泌物流出——这是尿道球腺（考珀腺）的液体，俗称"前液"。前液在排出时可能混有少量精子，这正是"体外射精并不保险"的原因之一。

睾丸因充血而胀大，大约能增加一半，并紧紧地贴向身体。精子沿着输精管向外侧出口运动，精囊与前列腺的分泌液汇入尿道。这一系列变化意味着射精已进入准备阶段：男性开始体验到一种"射精不可避免"的感觉——即使对阴茎的刺激完全停止，射精仍然要发生。这种迫切感大约持续 3 秒钟。

\subsubsection{高潮期}

射精分两步完成：先是尿道内口的括约肌收紧，精液被压入阴茎根部的尿道膨大处，此刻男性会感到一种压迫般的快感；随后尿道和环绕在阴茎根部的肌肉出现强烈收缩，把精液节律性地挤出体外。射出时收缩的力量可以很大，精液最远能射出约 28 厘米——当然，排出力度的个体差异极大，距离与快感强弱并不相关。

通常会发生 3—5 次强烈收缩，间隔约为 0.8 秒——与女性性高潮中阴道收缩的间隔完全一样。随着收缩的持续，力度渐渐变弱，间隔也失去规律。骨盆底部与大腿内侧的肌肉也会不自觉地收缩。男性性高潮的持续时间很少超过 15 秒，一般在 4—10 秒之间，比大部分女性的高潮略短。

\subsubsection{消退期}

高潮过后，男性的身体逐渐回到未唤起的状态。玛斯特斯和约翰逊夫妇观察到，男性的消退分两个阶段：第一阶段，阴茎仍保持半勃起，这种状态可以持续一段时间；第二阶段，勃起迅速而完全地消退，充血散去，睾丸降回原位，阴囊皮肤重新松弛下来。

与女性不同，男性在射精后有一段"不应期"——在这段时间里，无论受到多大的性刺激，都无法再次唤起勃起与射精。其机制与射精后催乳素升高、神经系统由兴奋转向抑制有关。不应期短的只有几分钟，长的可达数小时，并随年龄增长而延长，个体差异极大。本书其他章节已对不应期的原理与"二次性爱"的问题专门讨论过，此处不再展开。

年轻男性可以在尚未射精时退回平台期，让性紧张反复积累，甚至多次体验不伴随射精的高潮；但射精一旦发生，身体就会不可逆转地进入消退期。消退期里的男性普遍感到放松与困倦，而女性若在高潮后继续受到刺激，可以直接退回平台期、再度达到高潮。正因如此，性生活中的"事后时间"常常出现节奏差：男方只想睡，女方却期待温存与爱抚——把这段时间留给拥抱与交谈，对双方都是一种体贴。
